\begin{question}{22}{
    De Defensie Materieel Organisatie test de betrouwbaarheid van automatische hoogtemeters die gebruikt worden bij \enquote{High Altitude Low Opening} parachutesprongen.
    Deze hoogtemeters sturen op basis van luchtdruk een signaal naar het \enquote{Automatic Activation Device}, dat ervoor zorgt dat de parachute automatisch geopend wordt bij kritieke luchtdrukwaarden.
    
    De hoogte $X$ (in meters) waarop een willekeurige hoogtemeter uit de fabriek de kritieke luchtdruk meet, is normaal verdeeld met een gemiddelde van $\mu = \SI{1200}{\meter}$ en een standaardafwijking van $\sigma = \SI{35}{\meter}$.
}

    \subquestion{4}{
        De fabrikant van de hoogtemeters keurt de slechtst presterende \SI{2.5}{\percent} van de geproduceerde hoogtemeters direct af 
        (oftewel de hoogtemeters die een veel te lage hoogte registreren).
        Bereken de kritieke grenswaarde $g$ (in meters) waaronder een hoogtemeter wordt afgekeurd.
        Rond je antwoord af op één decimaal.
    }
    \solution{
        Er is gegeven dat de kansvariabele $X$ de hoogte meet waarop een willekeurige hoogtemeter uit de fabriek de kritieke luchtdruk meet, 
        en dat $X$ normaal verdeeld is met $\mu = 1200$ en $\sigma = 35$.

        We willen de grenswaarde $g$ berekenen waarvoor $X$ slechts met \SI{2.5}{\percent} kans, oftewel kans is $0.025$, een kleinere waarde aanneemt.
        Met andere woorden, we willen $g$ vinden met $P(X \le g) = 0.025$.\rubric{1}

        Dit kunnen we berekenen met de functie \enquote{InvNorm}:
        \begin{align*}
            g   &= \invnorm(\opp=0.025, \mu=1200, \sigma=35) \\
                &\approx 1131.40126.\rubric{2}
        \end{align*}

        Afgerond op één decimaal is de kritieke grenswaarde gelijk aan $g \approx 1131.4$. \rubric{1}
    }

    \subquestion{4}{
        De fabrikant slaagt erin om de hoogtemeters beter te kalibreren, waardoor de standaardafwijking onder operationele omstandigheden kan worden verkleind tot $\sigma=25$. 
        De gemiddelde hoogte $\mu = \SI{1200}{\meter}$ blijft wel gelijk.
        Een hoogtemeter weigert dienst als deze een hoogte meet die lager is dan de zojuist in subvraag (a) berekende grenswaarde $g$.
        Bereken de kans dat een willekeurige hoogtemeter dienst weigert.
    }
    \solution{
        De kansvariabele $X$ is in dit geval normaal verdeeld is met $\mu = 1200$ en $\sigma = 25$.\rubric{1}
        De hoogtemeter weigert dienst als de hoogte $X$ waarop de exacte luchtdruk wordt gemeten kleiner dan de berekende grenswaarde $1131.4$ is, oftewel $P(X \leq 1131.4)$.\rubric{1}
        Dit kunnen we berekenen met de functie \enquote{normalcdf}:
        \begin{align*}
            P(X \leq g \approx 1131.4)  &= \normalcdf(\lb=-\GRinfinity, \ub=1131.4, \mu=1200, \sigma=25) \\
                                        &\approx 0.0030.\rubric{2}
        \end{align*}
    }

    \subquestion{8}{
        De commandant wil een peloton van $n$ parachutisten tegelijk laten springen onder deze operationele omstandigheden.
        De hoogtemeters functioneren onafhankelijk van elkaar. 
        De commandant eist dat de kans dat er minstens één hoogtemeter dienst weigert binnen het peloton, hoogstens \SI{10}{\percent} mag zijn.
        Bepaal het maximale aantal parachutisten $n$ dat aan deze sprong mag deelnemen om nog aan de veiligheidseis van de commandant te voldoen.
    }
    \solution{
        Stel $Y$ is de kansvariabele die het aantal dienstweigerende hoogtemeters telt.
        Aangezien de hoogtemeters onafhankelijk van elkaar worden ingezet, is $Y$ een binomiale kansvariabele. \rubric{1}
        De parameters van deze binomiale verdeling zijn een onbekende $n$ en de succeskans $p \approx 0.0030$ (kans dat een hoogtemeter dienst weigert).\rubric{2}
        
        We willen nu de maximale waarde voor $n$ bepalen zodanig dat geldt $P( Y \geq 1) = 1 - P(Y \leq 0) \leq 0.10$. \rubric{1}

        Aangezien $n$ een geheel getal moet zijn, maken we met de grafische rekenmachine een tabel:
        \begin{align*}
            y_1 = 1 - \binomcdf(n=X, p=0.0030, k=0) \rubric{1}
        \end{align*}

        Dit geeft ons de volgende waarde voor $X$:
        \begin{align*}
            n = 35: &\qquad 1 - \binomcdf(n=35, p=0.0030, k=0) \approx 0.09982\ (\text{voldoet}) \\
            n = 36: &\qquad 1 - \binomcdf(n=36, p=0.0030, k=0) \approx 0.10252\ (\text{voldoet niet meer}) \rubric{2} 
        \end{align*}

        Het maximale aantal parachutisten dat aan deze sprong mag deelnemen om nog aan de veiligheidseis van de commandant te voldoen is dus $n = 35$.
    }

    \subquestion{6}{
        Stel dat de fabrikant de hoogtemeters nog verder doorontwikkelt, waardoor de kans op een dienstweigerende hoogtemeter nog verder daalt tot $0.001$.
        Bereken de verwachtingswaarde en de standaardafwijking van het aantal dienstweigerende hoogtemeters, gegeven het maximale aantal parachutisten zoals berekend bij subvraag (c).
    }
    \solution{
        In dat geval is het aantal dienstweigerende hoogtemeters $Y$ binomiaal verdeeld met parameters $n = 35$ en $p = 0.001$.\rubric{2}

        De formules voor de verwachtingswaarde $E(Y)$ en de standaardafwijking $\sigma(Y)$ zijn dan als volgt:
        \begin{align*}
            E(Y) &= n \cdot p = 35 \cdot 0.001 = 0.035. \rubric{2} \\
            \sigma(Y) &= \sqrt{n \cdot p \cdot (1 - p)} = \sqrt{35 \cdot 0.001 \cdot (1 - 0.001)} \approx 0.1870.\rubric{2}
        \end{align*}
    }
\end{question}